% =====================================================================
%  Relatório — Entrega 2 (Janela deslizante: GBN e SR) — RDT-UnB — TR2 2026/2
%  Limite: no máximo 8 páginas.
%  Os números, tabelas e figuras vêm de report/gerado/ e report/figuras/,
%  produzidos por:
%     python scripts/entrega02_executa.py --server 137.131.178.229 --group grupo07
%     python scripts/entrega02_analisa.py
%  Trechos a preencher estão marcados com \preencher{...} (vermelho no PDF).
% =====================================================================
\documentclass[11pt,a4paper]{article}

\usepackage[T1]{fontenc}
\usepackage[utf8]{inputenc}
\usepackage[brazil]{babel}
\usepackage[a4paper,margin=2.2cm]{geometry}
\usepackage{amsmath,amssymb}
\usepackage{graphicx}
\usepackage{booktabs}
\usepackage{array}
\usepackage{float}
\usepackage{xcolor}
\usepackage{tikz}
\usetikzlibrary{automata,positioning,arrows.meta}
\usepackage[nottoc,numbib]{tocbibind} % numera a seção de referências
\usepackage{hyperref}
\hypersetup{colorlinks=true,linkcolor=blue,citecolor=blue,urlcolor=blue}

\graphicspath{{figuras/}}
\addto\captionsbrazil{\renewcommand{\refname}{Referências bibliográficas}}
\setlength{\emergencystretch}{3em} % evita linhas estouradas por nomes em \texttt

% Marcador de conteúdo pendente (remova antes de entregar).
\newcommand{\preencher}[1]{\textcolor{red}{[#1]}}

% Figura que mostra um quadro de espaço reservado enquanto o arquivo não existe.
% Uso: \figura{arquivo.png}{0.9\linewidth}{Legenda}{fig:rotulo}
\newcommand{\figura}[4]{%
  \begin{figure}[H]
    \centering
    \IfFileExists{figuras/#1}%
      {\includegraphics[width=#2]{#1}}%
      {\fbox{\parbox[c][4cm][c]{#2}{\centering\textcolor{red}{\ttfamily figuras/\detokenize{#1}}}}}%
    \caption{#3}
    \label{#4}
  \end{figure}}

% Peças geradas por scripts/entrega02_analisa.py (com fallback se ainda não existem).
\newcommand{\gerado}[1]{\IfFileExists{gerado/#1}{\input{gerado/#1}}%
  {\preencher{rodar scripts/entrega02\_analisa.py para gerar gerado/#1}}}
\IfFileExists{gerado/dados.tex}{\newcommand{\eArquivo}{\texttt{small}}
\newcommand{\eSegmentos}{2\,048}
\newcommand{\eTimeout}{1\,000}
\newcommand{\eSeed}{42}
\newcommand{\eQuando}{26/10/2026, entre 15:00 e 15:42}
\newcommand{\eExecucoes}{21}
\newcommand{\eExecucoesOk}{10}
\newcommand{\eRttMediano}{18,68}
}{%
  \newcommand{\eArquivo}{\preencher{arquivo}}\newcommand{\eSegmentos}{\preencher{N}}%
  \newcommand{\eTimeout}{1000}\newcommand{\eSeed}{42}%
  \newcommand{\eQuando}{\preencher{data}}\newcommand{\eExecucoes}{21}%
  \newcommand{\eExecucoesOk}{\preencher{n}}\newcommand{\eRttMediano}{\preencher{RTT}}}

% Estilo das máquinas de estados.
\tikzset{
  fsm/.style={>=Stealth, auto, node distance=1.6cm and 3.4cm},
  st/.style={draw, rounded corners, rectangle, minimum width=2.3cm,
             minimum height=0.85cm, fill=blue!6, font=\small\ttfamily, align=center},
  lbl/.style={font=\scriptsize, align=center},
}

% ---------------------------------------------------------------------
\begin{document}

% ======================= CAPA ========================================
\begin{center}
  {\large\textbf{Universidade de Brasília --- CIC0236 Teleinformática e Redes 2 --- 2026/2}}\\[0.5em]
  {\Large\textbf{RDT-UnB: janela deslizante com Go-Back-N e Selective Repeat --- Relatório da Entrega 2}}\\[0.5em]
  Grupo \texttt{grupo07}: Lucas Adeodato de Oliveira (231013645), Lurian Correia Lima (222005395), Atila Sousa Fernandes (231013510)
\end{center}

% ======================= 1. INTRODUÇÃO ===============================
\section{Introdução}

Na Entrega~1 mostramos que o Stop-and-Wait (SAW) é limitado pelo RTT: com um
único segmento em trânsito, o cliente passa mais de 99\% do tempo esperando.
Nesta entrega implementamos dois protocolos de janela deslizante (\emph{pipelining})
sobre o mesmo protocolo RDT-UnB~v1.0 e o mesmo servidor \emph{pull-based}: o
Go-Back-N (GBN) e o Selective Repeat (SR)~\cite{kurose}. Nos dois, o cliente
mantém até $N$ \texttt{REQ}s enviados sem \texttt{DATA} de volta; eles diferem no
tratamento de perdas e de segmentos fora de ordem. Comparamos SAW, GBN e SR com
janelas de 4, 8 e 16, sem perda e com 5\% e 10\% de perda simulada.

% ======================= 2. IMPLEMENTAÇÃO ============================
\section{Implementação}

O cliente ganhou as opções \texttt{--protocol \{saw,gbn,sr\}}, \texttt{--window N},
\texttt{--loss P} e \texttt{--seed S}. Handshake (\texttt{HELLO}/\texttt{OK}),
encerramento (\texttt{FIN}/\texttt{DONE} com MD5) e o CSV de métricas são os
mesmos da Entrega~1. Só o laço de transferência muda
(\texttt{src/saw.py}, \texttt{src/gbn.py} e \texttt{src/sr.py}). Como o servidor
só envia \texttt{DATA} quando recebe \texttt{REQ}, a "janela" é a de pedidos: o
cliente faz o papel de emissor (decide o que pedir e quando reenviar) e também de
receptor (aceita, descarta ou guarda os \texttt{DATA}). O timeout é fixo, como no
SAW, e o \texttt{ACK|seq=n} é enviado assim que o \texttt{DATA} $n$ é aceito.

\subsection{Go-Back-N}

O cliente guarda \texttt{base} (segmento mais antigo sem \texttt{DATA}) e
\texttt{next} (próximo a pedir) e mantém a janela cheia, com
$\texttt{next} - \texttt{base} \le N$. O receptor só aceita o \texttt{DATA} da
\texttt{base}: qualquer outro (fora de ordem, atrasado ou duplicado) é
descartado. Há um único temporizador, o da \texttt{base}, contado a partir do
último \texttt{REQ} dela. Quando ele vence, o cliente reenvia \textbf{todos} os
\texttt{REQ}s de \texttt{base} até \texttt{next}$-1$ e soma 1 à coluna
\texttt{retransmissions} de cada um (Figura~\ref{fig:fsm-gbn}).

\begin{figure}[H]
\centering
\begin{tikzpicture}[fsm]
  \node[st] (hs) {HELLO\_SENT};
  \node[st, right=of hs, minimum width=3.6cm] (w) {JANELA\\{\scriptsize $[\mathtt{base},\,\mathtt{next})$}};
  \node[st, right=of w] (fin) {FIN\_SENT};
  \draw[->] (hs) -- node[lbl, above] {OK / base$\leftarrow$next$\leftarrow$0} (w);
  \draw[->] (w) -- node[lbl, above] {base $=$ total /\\MD5, envia FIN} (fin);
  \draw[->] (w) edge[loop above, looseness=4] node[lbl, align=left]
    {next $<$ base$+N$ / REQ next, next$+{+}$\\
     DATA seq$=$base / ACK, guarda, base$+{+}$} (w);
  \draw[->] (w) edge[loop below, looseness=4] node[lbl, align=left]
    {DATA seq$\neq$base / \textbf{descarta}\\
     timeout(base) / \textbf{reenvia REQ base \dots\ next$-1$}} (w);
\end{tikzpicture}
\caption{Máquina de estados do cliente GBN (rótulos \emph{evento / ação}).
Esgotadas 10 retransmissões da mesma base, a sessão termina com erro.}
\label{fig:fsm-gbn}
\end{figure}

\subsection{Selective Repeat}

O SR mantém a mesma janela de pedidos, mas cada segmento pendente tem o
\textbf{próprio} temporizador. Um \texttt{DATA} de qualquer segmento pendente é
aceito, confirmado com \texttt{ACK} e guardado num \emph{buffer} indexado por
\texttt{seq}, mesmo que chegue fora de ordem. A \texttt{base} desliza por cima
de todos os segmentos consecutivos já presentes no \emph{buffer}, liberando
espaço na janela para novos pedidos. Quando o temporizador de um segmento vence,
\textbf{só} o \texttt{REQ} daquele segmento é reenviado. \texttt{DATA} de
segmento que já não está pendente (duplicata) é descartado. Ao final, o arquivo
é remontado concatenando o \emph{buffer} em ordem de \texttt{seq}
(Figura~\ref{fig:fsm-sr}).

\begin{figure}[H]
\centering
\begin{tikzpicture}[fsm]
  \node[st] (hs) {HELLO\_SENT};
  \node[st, right=of hs, minimum width=3.6cm] (w) {JANELA\\{\scriptsize pendentes + buffer}};
  \node[st, right=of w] (fin) {FIN\_SENT};
  \draw[->] (hs) -- node[lbl, above] {OK / base$\leftarrow$next$\leftarrow$0} (w);
  \draw[->] (w) -- node[lbl, above] {base $=$ total / remonta,\\MD5, envia FIN} (fin);
  \draw[->] (w) edge[loop above, looseness=4] node[lbl, align=left]
    {next $<$ base$+N$ / REQ next, inicia timer(next), next$+{+}$\\
     DATA seq$=i$ pendente / ACK $i$, \textbf{buffer}[$i$], desliza base} (w);
  \draw[->] (w) edge[loop below, looseness=4] node[lbl, align=left]
    {DATA de seq já recebido / descarta\\
     timeout($i$) / \textbf{reenvia só REQ $i$}, reinicia timer($i$)} (w);
\end{tikzpicture}
\caption{Máquina de estados do cliente SR. Cada segmento pendente $i$ tem seu
próprio temporizador; até 10 retransmissões por segmento.}
\label{fig:fsm-sr}
\end{figure}

\subsection{Simulação de perda}

Escolhemos a abordagem (a): \textbf{descartar o \texttt{DATA} recebido com
probabilidade $P$} antes de processá-lo. O pacote chega do socket, mas o
cliente sorteia um número com \texttt{random.Random(seed)} (opção
\texttt{--seed}, padrão 42). Se o número for menor que $P$ (\texttt{--loss}), o
\texttt{DATA} é ignorado como se tivesse se perdido na rede, e o segmento
continua pendente até o timeout. Escolhemos descartar o \texttt{DATA}, e não o
\texttt{ACK}, porque no RDT-UnB quem dispara a retransmissão é a falta do
\texttt{DATA}: o servidor é \emph{pull-based} e não reenvia nada por conta
própria, e o próprio \texttt{ACK} não é processado pelo servidor, que
responde \texttt{ERROR|msg=tipo\_desconhecido\_ACK} a cada um. Assim, um
\texttt{ACK} omitido não teria efeito nenhum sobre a transferência,
enquanto um \texttt{DATA} descartado exercita exatamente o caminho de perda,
timeout e retransmissão de cada protocolo. A semente fixa torna os cenários
reproduzíveis.

No CSV, a coluna \texttt{loss\_simulated} (o campo \emph{simulated\_loss} do
enunciado) vale 1 se pelo menos um \texttt{DATA} daquele segmento foi descartado
pela simulação e 0 caso contrário. Isso separa a perda simulada das
retransmissões causadas por perda natural da rede: um segmento com
\texttt{retransmissions}$>0$ e \texttt{loss\_simulated}$=0$ foi reenviado por
perda real ou, no GBN, por ter sido arrastado pela retransmissão da janela. Na
demonstração ao vivo, o cliente roda com \texttt{--loss 0} (padrão).

% ======================= 3. METODOLOGIA ==============================
\section{Metodologia}

Os \eExecucoes{} cenários (SAW; GBN e SR com $N = 4$, 8 e 16; perdas de 0, 5 e
10\%) foram executados contra o servidor \texttt{137.131.178.229:8080} em
\eQuando, com o arquivo \eArquivo{} (\eSegmentos{} segmentos de 512\,B), timeout
fixo de \eTimeout\,ms e semente \eSeed. O script
\texttt{scripts/entrega02\_executa.py} roda os cenários em ordem intercalada
(para cada perda, todos os protocolos), para que uma piora momentânea da rede
não afete só um protocolo. O script \texttt{scripts/entrega02\_analisa.py} gera
as tabelas e os gráficos a partir dos CSVs. O throughput médio é
$\text{bytes recebidos} \times 8 / \text{tempo total}$, com o tempo medido do
primeiro \texttt{REQ} até o último \texttt{DATA}. As retransmissões são a soma
da coluna \texttt{retransmissions}.

% ======================= 4. RESULTADOS ===============================
\section{Resultados}

\subsection{Integridade}

\gerado{tabela_md5.tex}

\subsection{Tabela comparativa}

\gerado{tabela_comparativa.tex}

\subsection{Gráficos}

\figura{e2_janela_sem_perda.png}{0.72\linewidth}{Throughput médio sem perda em
função do tamanho da janela, comparado ao ideal $N \times L/\mathrm{RTT}$. Sem
perda, as curvas de GBN e SR praticamente se sobrepõem.}{fig:janela}

\figura{e2_segmentos_w16_p10.png}{0.92\linewidth}{GBN e SR com janela 16 e 10\%
de perda. Topo: throughput por segmento (coluna \texttt{throughput\_kbps},
média móvel de 16 segmentos); um segmento retransmitido inclui o timeout no seu
tempo. Abaixo: retransmissões por segmento. Os traços marcam os segmentos com
perda simulada.}{fig:segmentos}

% ======================= 5. ANÁLISE ==================================
\section{Análise}

\gerado{analise.tex}

% ======================= 6. CONCLUSÃO ================================
\section{Conclusão}

GBN e SR receberam o arquivo íntegro em todos os cenários
(\eExecucoesOk{}/\eExecucoes{} execuções com MD5 conferido e
\texttt{DONE|ok=1}), o que confirma a remontagem correta mesmo com segmentos
fora de ordem no SR. \gerado{conclusao.tex}

% ======================= REFERÊNCIAS =================================
\begin{thebibliography}{9}
  \bibitem{kurose}
    J. F. Kurose e K. W. Ross, \textit{Computer Networking: A Top-Down Approach},
    8th ed. Hoboken: Pearson, 2021.
  \bibitem{rfc768}
    J. Postel, ``User Datagram Protocol,'' RFC 768, IETF, 1980.
    Disponível em: \url{https://www.rfc-editor.org/rfc/rfc768}.
  \bibitem{rfc1321}
    R. Rivest, ``The MD5 Message-Digest Algorithm,'' RFC 1321, IETF, 1992.
    Disponível em: \url{https://www.rfc-editor.org/rfc/rfc1321}.
\end{thebibliography}

\end{document}
